\subsection*{2.1 The follow-up problem in literature-based discovery}\label{the-follow-up-problem-in-literature-based-discovery}

Literature-based discovery formalizes a simple observation: two facts can each be established in the
literature while the conclusion that joins them has never been written down. Swanson's original
formulation connects an entity \textbf{A} to an intermediate concept \textbf{B}, and \textbf{B} to an outcome \textbf{C},
where the \textbf{A--C} link is absent \cite{swanson1986}. In \emph{open} discovery one fixes \textbf{A} and searches for
reachable \textbf{C}; in \emph{closed} discovery one fixes both \textbf{A} and \textbf{C} and asks whether a plausible \textbf{B}
bridges them. Four decades of work have refined how the intermediate terms are linked, filtered, and
ranked, and have applied the paradigm across biomedicine --- including our own use of an ABC pipeline to
connect a plasma-metabolite signal to a druggable target in cardiac arrest \cite{henry2021}.

The paradigm's persistent weakness is not generation but \emph{triage}. An ABC search over a large corpus
returns thousands of candidate links, and the ranking signals available to classical LBD are properties
of the literature --- how rare a co-mention is, how strong an intermediate association looks. Rarity in the
literature, however, tracks \emph{obscurity} at least as strongly as it tracks \emph{importance}: a gene can look
novel simply because no one has studied it. The result is a long ranked list that no laboratory can work
through, and the field has said so of itself. Reviewing a proposal submitted for funding consideration,
three reviewers named the problem directly --- LBD \emph{``generates an enormous number of hypotheses, almost none
of which ever get followed up''}. Standard evaluations of LBD (time-slicing, link prediction) measure
whether the generator would have \emph{rediscovered} known links, not whether a specific proposed hypothesis is
supported by independent experimental data it has not seen \cite{henry2021}. We are precise about the scope of the answer this enables: adjudicating a
hypothesis against a held experimental substrate closes the \emph{triage} loop --- which of the machine's
hypotheses is worth pursuing --- not the experimental-follow-up loop itself. What has been missing is that
adjudication step: a way to put each machine-generated hypothesis to a data test and record a verdict ---
including a \emph{no}.

\subsection*{2.2 A held experimental substrate, and what counts as a receipt at each hop}\label{a-held-experimental-substrate-and-what-counts-as-a-receipt-at-each-hop}

The adjudication step needs data that exists independently of the hypotheses being posed. We use a
genome-scale CRISPRi Perturb-seq resource in primary human CD4+ T cells \cite{zhu2025},
in which thousands of genes are individually knocked down and single-cell RNA sequencing reads out the
transcriptional consequence of each perturbation across stimulation conditions. Because the measurements
were made before --- and without reference to --- any triple we test, they function as a \emph{held substrate}: the
referee queries them retrospectively rather than commissioning new experiments. The analysis is
CPU-feasible because it runs over the study's aggregated supplementary tables --- per-guide knockdown
efficiency, differential-expression statistics, a Th1/Th2 polarization signature, and a
cluster-to-autoimmune-disease enrichment --- rather than the raw multi-million-cell matrices. This substrate
is not a universal oracle of biological truth; it is a bounded, independent adjudication surface for
mechanistic T-cell claims --- valuable precisely because, for a hypothesis it \emph{can} speak to, it returns
support, refutation, or \emph{untested} with a receipt, and stays silent (rather than guessing) on the rest.

Crucially, the four tables do not all supply the same \emph{kind} of evidence, and we keep the distinction
explicit throughout. The first three hops --- did the knockdown work, did the perturbation produce a real
transcriptional effect, did that effect shift the T-cell program --- are backed by \textbf{experimental
receipts}: direct measurements from the perturbation data. The fourth hop --- does the perturbed gene's
downstream signature enrich for a disease --- is an \textbf{association receipt}: the disease labels derive from
genetic evidence (GWAS and related association data, without colocalization or linkage-disequilibrium
control), so a positive verdict at this hop is a \emph{nomination}, not a claim of experimental disease
causality. Reading the substrate as a direct experimental readout for the first three hops but an
association signal for the fourth is what keeps the eventual claims calibrated.

\subsection*{2.3 Agentic language models as a scientific instrument}\label{agentic-language-models-as-a-scientific-instrument}

The generation, adjudication, and provenance in this work were carried out inside an \emph{agentic} language-
model environment, and because that setting is newer to the discovery-informatics reader than the biology
or the LBD, we describe what it does and --- more importantly --- what discipline makes it trustworthy for
science. Contemporary language models can be given \emph{tools}: they can call code, query a web API, or read a
file, and act on the returned values rather than on their own recollection. They can also be composed into
an \emph{actor--critic} configuration, in which one model performs an analysis and a second, independent model
reviews it.

Two design commitments make this usable as an instrument rather than a source of confident narration.
First, we separate deterministic \emph{lookups} from \emph{judgment}. Every data receipt in Wayfinder --- a knockdown
efficiency, an effect size, an enrichment odds ratio --- is produced by deterministic code, reproducible and
free of model discretion; the language model \emph{interprets} receipts and assembles the verdict, but it never
asserts a biological fact that a table value does not support. Visible agency is reserved for judgment, not
for data retrieval. Second, the work runs on a scientific workbench --- Claude Science, a persistent-kernel
environment --- in which an author model and an independent reviewer model operate at separate checkpoints:
the reviewer verifies each reported number against the underlying artifacts and enforces calibrated
language on the output, and in this work it flagged and removed overstated words (for example ``validated''
and ``definitive'') from the platform's own text. We are precise about what this independence is and is not:
it is \emph{role, model, and checkpoint} independence --- a distinct reviewer model at distinct verification
points --- within a single model family, not cross-vendor independence. Independence across model families
is provided separately and externally (Section 4.6), by replicating the finding with a different vendor's
models. With that boundary stated, the agentic workbench lets a single researcher run generation,
data-grounded adjudication, and a self-audit of the result's language on one bench --- the capability the
rest of this paper puts to work.
